\par\noindent\textbf{(1)} 对长度dx的小段应用Kirchhoff定律得到 ∂\_xV=−R′I−L′∂\_tI 与 ∂\_xI=−G′V−C′∂\_tV。再对第一式求x导并代入第二式，得到 ∂\_x²V=L′C′V\_tt+(R′C′+L′G′)V\_t+R′G′V；电流满足同型方程。
\par\noindent\textbf{(2)} 令V∝e\textasciicircum{}\{−γx−iωt\}，则γ²=(R′−iωL′)(G′−iωC′)，取 Reγ≥0 保证被动媒质向+x不自发增长。由−∂\_xV=(R′−iωL′)I 得 Z₀=V/I=√[(R′−iωL′)/(G′−iωC′)]，选择使功率向正方向输送的根支。
\par\noindent\textbf{(3)} 若R′≪ωL′、G′≪ωC′，展开 γ≈α−iβ，得到 β≈ω√(L′C′)、α≈(R′/2)√(C′/L′)+(G′/2)√(L′/C′)。于是 v\_p≈1/√(L′C′)，α≥0且两种损耗以串、并联方式相加，反映线路的高频能量耗散。
\par\noindent\textbf{(4)} 在RC线 L′→0、G′→0 时电报方程退化为 ∂\_x²V=R′C′V\_t，即 V\_t=D\_RC V\_xx，D\_RC=1/(R′C′)。传播常数 γ=√(−iωR′C′)=(1−i)√(ωR′C′/2)，相位与衰减长度相同量级，低频信号表现扩散而非无损波。
